% GENERATED FILE. Edit canonical lessons, then run npm run generate:books.
% canonical-source-hash: fnv1a64:9b33e9153ed908e6
% canonical-lessons: HI-C114-mafi, HI-C114-koibaat, HI-C114-badhai, HI-C114-shubhkamnaen, HI-C114-nimantran

\chapter{Sorry, No Problem, Congratulations}
\label{ch:hi-sorry-no-problem}
\hlchaptermodality{114}

\begin{chapteropening}
\textbf{By the end of this chapter:} \emph{I can apologise, answer an apology, congratulate, give good wishes, and talk about an invitation.}

Complete HI-C114-nimantran: the courtesies that come around a visit.
\end{chapteropening}

\section[māfī]{\hi{माफ़ी} (māfī) — an apology, pardon}

\label{lesson:HI-C114-mafi}



\begin{quote}
Before the new one: say the Hindi for that very one, then the Hindi for maybe.
\end{quote}

\subsection*{You'll want to know: \hi{माफ़ी}}
\textbf{\hi{माफ़ी}} (\emph{māfī}) — \textquotedblleft{}an apology, pardon\textquotedblright{}, the noun.

\textbf{\hi{माफ़ी} \hi{चाहता} \hi{हूँ}} — \textquotedblleft{}I beg your pardon\textquotedblright{} (a woman says \textbf{\hi{चाहती}}). It shares its root with \textbf{\hi{माफ़} \hi{कीजिए}}.

\begin{grammarlens}[title={What you've built}]
The apology as a thing you ask for.
\end{grammarlens}

\subsection*{Your turn}
\begin{itemize}
\raggedright
  \item \emph{Say it:} \emph{māfī}
  \item \emph{Say it:} \emph{māfī}, once more
  \item \emph{Say it:} say \emph{māfī cāhtā hūṁ}
  \item \emph{Recall:} say the Hindi for that very one, then the Hindi for maybe, then say \emph{māfī} again
\end{itemize}

\begin{tcolorbox}[breakable,colback=teal!4,colframe=teal!35!black,title={Before you move on}]
What does \hi{माफ़ी} mean? (\textquotedblleft{}An apology, pardon\textquotedblright{}.) Say it once more.
\end{tcolorbox}
\section[koī bāt nahīṁ]{\hi{कोई} \hi{बात} \hi{नहीं} (koī bāt nahīṁ) — no problem, never mind}

\label{lesson:HI-C114-koibaat}



\begin{quote}
Before the new one: say the Hindi for maybe, then the Hindi for an apology.
\end{quote}

\subsection*{You'll want to know: \hi{कोई} \hi{बात} \hi{नहीं}}
\textbf{\hi{कोई} \hi{बात} \hi{नहीं}} (\emph{koī bāt nahīṁ}) — \textquotedblleft{}no problem\textquotedblright{}, \textquotedblleft{}never mind\textquotedblright{}, literally \textquotedblleft{}no matter\textquotedblright{}.

It is the answer to \textbf{\hi{माफ़} \hi{कीजिए}}, and to thanks for a small favour.

\begin{grammarlens}[title={What you've built}]
The answer to an apology.
\end{grammarlens}

\subsection*{Your turn}
\begin{itemize}
\raggedright
  \item \emph{Say it:} \emph{koī bāt nahīṁ}
  \item \emph{Say it:} \emph{koī bāt nahīṁ}, once more
  \item \emph{Say it:} answer \emph{māf kījie} with \emph{koī bāt nahīṁ}
  \item \emph{Recall:} say the Hindi for maybe, then the Hindi for an apology, then say \emph{koī bāt nahīṁ} again
\end{itemize}

\begin{tcolorbox}[breakable,colback=teal!4,colframe=teal!35!black,title={Before you move on}]
What does \hi{कोई} \hi{बात} \hi{नहीं} mean? (\textquotedblleft{}No problem, never mind\textquotedblright{}.) Say it once more.
\end{tcolorbox}
\section[badhāī]{\hi{बधाई} (badhāī) — congratulations}

\label{lesson:HI-C114-badhai}



\begin{quote}
Before the new one: say the Hindi for an apology, then the Hindi for no problem.
\end{quote}

\subsection*{You'll want to know: \hi{बधाई}}
\textbf{\hi{बधाई}} (\emph{badhāī}) — \textquotedblleft{}congratulations\textquotedblright{}.

\textbf{\hi{बधाई} \hi{हो}!} — \textquotedblleft{}congratulations!\textquotedblright{}, for an exam passed, a new job, a birth.

\begin{grammarlens}[title={What you've built}]
Praise for good news.
\end{grammarlens}

\subsection*{Your turn}
\begin{itemize}
\raggedright
  \item \emph{Say it:} \emph{badhāī}
  \item \emph{Say it:} \emph{badhāī}, once more
  \item \emph{Say it:} say \emph{badhāī ho!}
  \item \emph{Recall:} say the Hindi for an apology, then the Hindi for no problem, then say \emph{badhāī} again
\end{itemize}

\begin{tcolorbox}[breakable,colback=teal!4,colframe=teal!35!black,title={Before you move on}]
What does \hi{बधाई} mean? (\textquotedblleft{}Congratulations\textquotedblright{}.) Say it once more.
\end{tcolorbox}
\section[śubhkāmnāẽ]{\hi{शुभकामनाएँ} (śubhkāmnāẽ) — good wishes}

\label{lesson:HI-C114-shubhkamnaen}



\begin{quote}
Before the new one: say the Hindi for no problem, then the Hindi for congratulations.
\end{quote}

\subsection*{You'll want to know: \hi{शुभकामनाएँ}}
\textbf{\hi{शुभकामनाएँ}} (\emph{śubhkāmnāẽ}) — \textquotedblleft{}good wishes\textquotedblright{}.

It opens on the \textbf{\hi{शुभ}} of \textbf{\hi{शुभ} \hi{रात्रि}}, \textquotedblleft{}auspicious\textquotedblright{}. \textbf{\hi{आपको} \hi{शुभकामनाएँ}} — \textquotedblleft{}good wishes to you\textquotedblright{}.

\begin{grammarlens}[title={What you've built}]
Wishes for what is coming, beside congratulations for what is done.
\end{grammarlens}

\subsection*{Your turn}
\begin{itemize}
\raggedright
  \item \emph{Say it:} \emph{śubhkāmnāẽ}
  \item \emph{Say it:} \emph{śubhkāmnāẽ}, once more
  \item \emph{Say it:} say \emph{āpko śubhkāmnāẽ}
  \item \emph{Recall:} say the Hindi for no problem, then the Hindi for congratulations, then say \emph{śubhkāmnāẽ} again
\end{itemize}

\begin{tcolorbox}[breakable,colback=teal!4,colframe=teal!35!black,title={Before you move on}]
What does \hi{शुभकामनाएँ} mean? (\textquotedblleft{}Good wishes\textquotedblright{}.) Say it once more.
\end{tcolorbox}
\section[nimantraṇ]{\hi{निमंत्रण} (nimantraṇ) — an invitation}

\label{lesson:HI-C114-nimantran}



\begin{quote}
Before the new one: say the Hindi for congratulations, then the Hindi for good wishes.
\end{quote}

\subsection*{You'll want to know: \hi{निमंत्रण}}
\textbf{\hi{निमंत्रण}} (\emph{nimantraṇ}) — \textquotedblleft{}an invitation\textquotedblright{}, the formal word.

\textbf{\hi{निमंत्रण} \hi{के} \hi{लिए} \hi{धन्यवाद}} — \textquotedblleft{}thank you for the invitation\textquotedblright{}, with \textbf{\hi{के} \hi{लिए}} and \textbf{\hi{धन्यवाद}}.

A written invitation has the three parts of the short message you learned to write: the line that names the reader, \textbf{\hi{प्रिय} \hi{मीरा},}; the sentences under it; and the writer's name at the foot.

\begin{grammarlens}[title={What you've built}]
Sorry, no problem, congratulations, wishes, and thanks for an invitation.
\end{grammarlens}

\subsection*{Your turn}
\begin{itemize}
\raggedright
  \item \emph{Say it:} \emph{nimantraṇ}
  \item \emph{Say it:} \emph{nimantraṇ}, once more
  \item \emph{Say it:} say \emph{nimantraṇ ke lie dhanyavād}
  \item \emph{Recall:} say the Hindi for congratulations, then the Hindi for good wishes, then say \emph{nimantraṇ} again
\end{itemize}

\begin{tcolorbox}[breakable,colback=teal!4,colframe=teal!35!black,title={Before you move on}]
What does \hi{निमंत्रण} mean? (\textquotedblleft{}An invitation\textquotedblright{}.) Say it once more.
\end{tcolorbox}
